\scanpage{160}
On ne suppose plus donné de $\underline{L}$

A'') Équivalence entre les données suivantes

(i) Une Section $\xi$ de faisceau \unsure{inversible|convenable} sur $S$, \inserted{qui est \uline{morphique}} et un isom
\[
\varphi:P_{\unsure{\xi|S}}^{N}\xrightarrow{\sim}P_{\xi}
\]

(ii) Une Donnée d'un morphisme
\[
j:X\longrightarrow P_{S}^{N}
\]
tel que \deleted{\unsure{$j^{*}\underline{O}$}}
\[
g_{*}(\underline{O}(1))\longrightarrow f_{*}(j^{*}(\underline{O}(1))
\]
soit un isomorphisme
% slanted
\margin{\uline{morphisme spécial} dans un projectif}

\margin{\unsure{Ou bien|Plan bis|Attention}, $2$}
\diagram{\begin{tikzcd}[ampersand replacement=\&,row sep=small]
{\underline{L}\ \text{morphique }\inserted{\text{spécial}}\text{ sur }X/S
  \;+\;\text{trivialisation projective}}\arrow[d]\\
{\text{Section }\xi\text{ sur }S,+\text{ \unsure{trivialisation} projective}}\arrow[d,equal]\\
{\text{Morphisme spécial }X\longrightarrow P^{N}_{S}}\arrow[d]\\
{\underline{L}'\text{ morphique spécial sur }X,\text{ et trivialisation projective}}
\end{tikzcd}}
\margin{\#\quad isom \deleted{\unsure{avec}} \unsure{n'est pas} \unsure{forcément} identique}

A''') Équivalence entre les données \deleted{\unsure{suivantes}}

(i) Une Section \deleted{\unsure{morphique et spéciale}} de faisceaux \inserted{\unsure{``morphique et spécial''}} sur $X$
% slanted
\margin{spécial = \unsure{$1^{\circ}$} coh. plt en dim $0$, et tel qu'il définisse un morphisme dans un projectif \unsure{i.e.} image directe loc. libre de \deleted{\unsure{genre positif}} \inserted{dim.} $N$}

(ii) Un fibré de \uline{Brauer-Severi} \deleted{\unsure{$X$}} sur $S$, associé à un \uline{fibré principal} $\unsure{\mathscr{P}|\mathfrak{P}}$ de groupe $\mathrm{GP}\deleted{\unsure{L}}(N)$

$+$ un morphisme \uline{spécial} $j:X_{\mathscr{P}}\longrightarrow P^{N}_{\mathscr{P}\deleted{\unsure{(E)}}}$

\deleted{\unsure{la sphère|la fibre}}

ce morphisme étant \uline{covariant} pour $\mathrm{GP}(N)$
% slanted
\margin{\deleted{\unsure{pour ces principaux} \unsure{fibrés}, \illegible{2} \unsure{introduire les|introduction} \unsure{conditions de passage} \unsure{l'invariance} \illegible{2} \unsure{sur les points} \unsure{où} $j^{*}\underline{O}(1)$ est \unsure{de} \illegible{1}} \unsure{au fond} \illegible{2} \unsure{quoi ?}}

Si on introduit \deleted{\unsure{$J$|$\mathfrak{J}$}} $\underline{\underline{\operatorname{Hom}}}_{S}(X,P^{N})$ où $\mathrm{GP}(N)$ opère, le sous-schéma \deleted{\illegible{1}} $J=\underline{\underline{\operatorname{Hom}}}_{S}^{\mathrm{spéc}}(X,P^N)$ [\unsure{si} $X$ \deleted{\unsure{propre}} projectif sur $S$ loc. \unsure{noeth}], alors \deleted{\unsure{pour les structures}} on voit que $\mathrm{GP}(N)$ opère \uline{sans pts fixes} sur $\underline{\underline{\operatorname{Hom}}}_{S}^{\mathrm{spéc}}(X,P^{N})$, et les structures A''' sont
% archive sheet number: 159
